\section{Problem Statement}

\begin{frame}{The Challenge}
  \begin{columns}[T,onlytextwidth]
    \column{.58\textwidth}
      \footnotesize
      \begin{itemize}
        \item AI-generated imagery is now visually indistinguishable from real photographs.
        \item Two entangled problems: \textbf{Authenticity} (real or fake?) and \textbf{Provenance} (re-transmitted or re-digitized?).
        \item Real-world detectors must survive this ``laundering'' pipeline of re-compression and re-digitization.
        \item We frame this as a \textbf{joint} problem: processing history changes what artifacts remain for authenticity.
      \end{itemize}

    \column{.38\textwidth}
      \centering
      \begin{tikzpicture}[
          node distance=5mm,
          box/.style={cookie node,minimum width=2.7cm,minimum height=.75cm,font=\scriptsize\bfseries,align=center},
          accentbox/.style={cookie accent node,minimum width=2.7cm,minimum height=.75cm,font=\scriptsize\bfseries,align=center}
        ]
        \node[box] (orig) {Original image};
        \node[box,below=of orig] (transfer) {Transfer /\\compression};
        \node[box,below=of transfer] (redigital) {Re-digitization};
        \node[accentbox,below=of redigital] (detector) {Detector};
        \draw[cookie arrow] (orig) -- (transfer);
        \draw[cookie arrow] (transfer) -- (redigital);
        \draw[cookie arrow] (redigital) -- (detector);
      \end{tikzpicture}
  \end{columns}
\end{frame}
